\documentclass[11pt]{article}

\usepackage{amsmath,amssymb}
\usepackage{xurl}
\usepackage[hidelinks]{hyperref}
\setlength{\emergencystretch}{2em}

\input{command}

% Local notation for this note.
\DeclareMathOperator{\poly}{poly}
\newcommand{\Id}{\mathrm{Id}}
\renewcommand{\ghissue}[1]{%
  \href{https://github.com/Dengnifer/MIPStarRE-A/issues/#1}{GitHub issue \##1}}

\title{The Error Term of the Combined-Lines Consistency Lemma}
\author{}
\date{2026-08-30}

\begin{document}
\maketitle

\section{At a glance}

\begin{itemize}
  \item \textbf{Current paired-line construction status:} for
  \texttt{lem:qld-xz-lines}, actual conditional pasting and restoration to
  the original question law are now proved, retaining the supplied point
  error. Polynomial specialization and the actual paired-line witness,
  including every opposite placement and axis-degree support, are now proved.
  The three-claim overlap estimate and its scalar polynomial absorption are
  also proved. The actual extended-line POVM, degree support, and averaged
  placed-overlap comparison are proved, as is the first-player quantitative
  estimate. The second-player X-removal and Z-deficit estimates and aggregate
  consumer sufficiency remain open. The printed existence theorem is not disproved; no final correction
  of the extended-line theorem is adopted by these results.
  \item \textbf{Difficulty: substantial proof and transport obligations.}
  The current verdict is not converged. The ordinary ten-session
  investigation ends with the unresolved obligations recorded below.
  The obstruction is that the error function claimed in the
  statement of the combined-lines lemma of the Pauli basis analysis,
  $\poly(m^2\eps, md/q)$, is not delivered by either of the two
  derivations printed in its source proof: one route establishes
  $m \cdot \poly(\eps, md/q)$, the other would give
  $m^2 \cdot \poly(\eps, md/q)$ and rests in addition on a distributional
  decomposition that the stated sub-line lemma does not supply.  The
  candidate correction is the first-route error form, subject to the full
  two-player construction and consumer sufficiency checks.
  \item \textbf{Key Mathlib inputs:} the
  Cauchy--Schwarz inequality for state-dependent operator distances
  (project-local, blueprint Chapter 12), mixture decompositions of
  question distributions, and elementary averaging over restricted
  line distributions; no new Mathlib surface beyond finite
  expectations and inner-product estimates already used by the
  low individual degree development.
  \item \textbf{Estimated weight:} two second-player estimates and witness
  assembly, followed by source and consumer transport; the aggregate remaining
  weight is not established. The remaining work comprises the second-player first-route estimates,
  assembly of both consistency fields, the explicit global-pair parameter
  absorption, and the source-question and field-answer transport obligations.
  Compiling an immediate wrapper does not establish these uses.
  \item \textbf{Mathlib/project split:} essentially \textbf{100\%
  project-local}.  The estimates use only the project's state-dependent
  distance calculus, the sub-line distribution, and elementary
  probability on finite sets.
  \item \textbf{Key inputs:} the combined point and line measurements and
  the sub-line distribution of the Pauli basis analysis in
  \cite{Ji2020MIPStar}, and the error function of the classical soundness
  theorem imported from \cite{Ji2020LowIndividualDegree}, into which the
  corrected error form must be absorbed.  The blueprint correction is tracked
  in \ghissue{173}.
\end{itemize}

\section{Key theorem forms}

Throughout, $(q,m,d)$ is an admissible parameter tuple of the Pauli basis
test, $\ket{\hat\psi}$ is the expanded state of the projective strategy
under analysis, succeeding with probability at least $1-\eps$, and all
approximations are with respect to $\ket{\hat\psi}$.  For lines
$\ell_X, \ell_Z$ in $\F_q^m$, the POVM $\{T^{\ell_X,\ell_Z}_{f_X,f_Z}\}$,
with outcomes pairs of polynomials of degree at most $md$ on the two
lines, is the two-basis combined line measurement, and
$\{\hat Q^{x,z}_{a,b}\}$ is the combined point measurement; their
consistency errors are $\delta_P = \poly(\eps, md/q)$ and
$\delta_Q = \poly(\eps)$ respectively, and
$\delta_{\mathrm{Line}} = \poly(\eps)$ is the line--point consistency
error of the strategy.  The lemma under discussion asserts that a
measurement $\{\hat Q^{\ell}_f\}$, built for each line $\ell$ in
$\F_q^{2m+2}$, is $\delta_{\mathrm{combine}}$-consistent, on average over
the line-point distribution over $\F_q^{2m+2}$, with the combined point
measurement on the extended coordinates, together with the
register-swapped counterpart.%
\footnote{In the mirror of the source consulted here,
\path{references/qpbt-paper/14_analysis_of_the_pauli_basis_test.tex}: the
lemma statement, with the error form $\poly(m^2\eps, md/q)$, is at
lines~1020--1034 (label \path{lem:qld-4-13}); the central display with
the first-route bound at lines~1134--1137
(label \path{eq:qld-combined-lines-consistency}); the three claims at
lines~1140--1239 (labels \path{claim:17-1}--\path{claim:17-3}); the
second-route conclusion $O(m^2\delta_P)$ at lines~1241--1245, preceded by
a duplicated sentence whose second copy carries a malformed reference;
the product-distribution estimate invoked there at lines~1055--1061
(label \path{eq:qld-xz-lines-restricted}); and the sub-line lemma at
lines~1063--1116 (label \path{lem:qld-sublines}).}
The comparison involves four forms of the error function
$\delta_{\mathrm{combine}}(\eps,m,d,q)$.
\begin{enumerate}
  \item \textbf{Source form.}
  $\delta_{\mathrm{combine}} = \poly(m^2\eps,\, md/q)$.  This is the form
  printed in the statement.  Neither printed derivation supports it.
  \item \textbf{Established form (first route).}
  $\delta_{\mathrm{combine}}
  = O\big(m\,(\eps^{1/4} + \delta_P^{1/4} + \delta_Q^{1/4}
  + \delta_{\mathrm{Line}}^{1/2})\big)$, a function of the form
  $m \cdot \poly(\eps,\, md/q)$.
  \item \textbf{Second-route form (conditional).}
  $\delta_{\mathrm{combine}} = O(m^2 \delta_P)
  = m^2 \cdot \poly(\eps,\, md/q)$, conditional on a joint-law
  decomposition that the sub-line lemma, as stated, does not provide.
  \item \textbf{Consumer requirement.}  The unique consumer, the
  polynomial-pair lemma, feeds $\delta_{\mathrm{combine}}$ into an error
  of the form
  $a(md)^a(\delta_{\mathrm{combine}}^b + q^{-b} + 2^{-bmd})$ imported
  from \cite{Ji2020LowIndividualDegree}; any bound
  $m^{O(1)} \cdot \poly(\eps, md/q)$ is absorbed by enlarging $a$ and
  shrinking $b$, so forms~2 and~3 both suffice downstream.
\end{enumerate}

\section{The source assertion and the two routes}

The source proof in \cite{Ji2020MIPStar} constructs $\hat Q^{\ell}_f$ by
averaging, over a distribution $D$ on tuples $(\ell, \ell_X, \ell_Z)$
provided by its sub-line lemma, the pushforward of
$T^{\ell_X,\ell_Z}$ along a combining map, and reduces the lemma to a
bound on the scalar consistency defect
\begin{equation}
  \Lambda \,=\,
  \E_{(\ell,\ell_X,\ell_Z) \sim D}\
  \E_{(x,z,\alpha,\beta) \in \ell}
  \sum_{f_X,f_Z} \bra{\hat\psi}
  \big(T^{\ell_X,\ell_Z}_{f_X,f_Z}\big)_{\mathsf{A}\mathsf{A}'} \ot
  \big(\Id - \hat{Q}^{x,z}_{f_X(x),f_Z(z)}\big)_{\mathsf{B}\mathsf{A}''}
  \ket{\hat\psi}\,,
  \tag{\path{eq:qld-combined-lines-consistency}}
\end{equation}
in which the point $(x,z,\alpha,\beta)$ is uniform on $\ell$ and the sum
ranges over pairs of line polynomials.  The proof then bounds $\Lambda$
twice, by different routes and with different results, and the printed
record is internally inconsistent: the display introducing $\Lambda$
asserts the first-route bound, while the closing lines assert the
second-route bound as the operative one; the lemma statement claims a
third form matching neither.

\paragraph{The first route.}
Three claims are chained.  The combined point measurement
$\hat Q^{x,z}$ is replaced by the ordered product of the two point
measurements at cost $\delta_Q^{1/2}$; the $X$-point factor is removed at
cost $m\,\delta_{\mathrm{Line}}^{1/2}$; and the remaining $Z$-point
correlation is shown to be within
$\sqrt{m}\,(\delta_P^{1/4} + \delta_Q^{1/4} + \eps^{1/4})$ of one.  The
factors of $m$ enter through the mixture structure of the sub-line
distribution: each restricted line distribution carries mixture weight
$1/(2m)$ in the line-point distribution, so an average bounded by
$\delta$ over the full distribution is bounded only by $2m\delta$ over a
restricted component and by $4m^2\delta$ over a product of two
components, and the Cauchy--Schwarz steps halve the exponents.  Chaining
the claims gives
\begin{equation*}
  \Lambda \,=\, O\big(m\,(\eps^{1/4} + \delta_P^{1/4} + \delta_Q^{1/4}
  + \delta_{\mathrm{Line}}^{1/2})\big)
  \,=\, m \cdot \poly(\eps,\, md/q)\,.
\end{equation*}
One intermediate assertion needs a harmless correction: the final line
of the proof of the second claim asserts the bound
$O(m^2 \delta_{\mathrm{Line}})$ for a consistency term that is close to
one, rather than for its deficit.  Read as a bound on the deficit ---
which is what the surrounding Cauchy--Schwarz argument requires and what
the invoked line--point consistency provides after the mixture
inflation --- the claim and its use are correct, with the square root of
the deficit giving the stated cost $m\,\delta_{\mathrm{Line}}^{1/2}$.

\paragraph{The second route.}
After the chain, the source restarts: it asserts that, by Property~2 of
its sub-line lemma, $\Lambda$ is a mixture, over types
$\mathrm{v}_1, \mathrm{v}_2 \in \{\mathrm{ALine}, \mathrm{DLine}\}$ and
indices $i, j \in \{1, \dots, m\}$, of terms
\begin{equation*}
  \E_{(\ell_X, x) \sim D_{\mathrm{v}_1,i}}\
  \E_{(\ell_Z, z) \sim D_{\mathrm{v}_2,j}}
  \sum_{f_X,f_Z} \bra{\hat\psi}
  \big(T^{\ell_X,\ell_Z}_{f_X,f_Z}\big)_{\mathsf{A}\mathsf{A}'} \ot
  \big(\Id - \hat{Q}^{x,z}_{f_X(x),f_Z(z)}\big)_{\mathsf{B}\mathsf{A}''}
  \ket{\hat\psi}
\end{equation*}
over two \emph{independently} sampled restricted line-point pairs,
bounds each term by $O(m^2\delta_P)$ using the restricted-distribution
consequence of the two-basis line lemma, and concludes
$\Lambda = O(m^2\delta_P)$.  The invoked decomposition concerns the
\emph{joint} law of the pair $((\ell_X, x), (\ell_Z, z))$, where $x$ and
$z$ are the two blocks of the single point sampled on $\ell$.  Property~2
of the sub-line lemma does not assert this joint law: it asserts the
component-wise product law of the pair of \emph{lines}
$(\ell_X, \ell_Z)$, and the conditional uniformity of each point on its
line separately --- equivalently, the laws of the two pairs
$((\ell_X, x), (\ell_Z, z'))$ and $((\ell_X, x'), (\ell_Z, z))$ in which
the \emph{other} point is resampled independently.  These marginal
statements are exactly what the first route's claims consume, each
involving one point at a time; the joint decomposition used by the
second route is a strictly stronger property and is not established by
the lemma it is attributed to.

\section{Why neither route yields the stated form}

The stated error form is $\poly(m^2\eps, md/q)$, a function of the two
quantities $m^2\eps$ and $md/q$ that vanishes as both vanish.  The
first-route bound is not dominated by any such function: along the
scaling $\eps = m^{-4}$ with $q/(md) \to \infty$, both $m^2\eps = m^{-2}$
and $md/q$ tend to $0$, while the term $m\,\eps^{1/4} = 1$ does not.
Hence the established derivation supports no bound of the claimed form.
The second-route bound, even granting the missing decomposition, carries
the separate prefactor $m^2$: writing $\delta_P = \poly(\eps, md/q)$
with some fixed positive exponents, the contribution
$m^2 (md/q)^{c}$ is not dominated by any vanishing function of $m^2\eps$
and $md/q$, as the scaling $\eps = 0$, $md/q = m^{-2/c}$ shows, and the
$\eps$-contribution fails similarly for sublinear exponents.  Neither
observation asserts that the \emph{conclusion} of the lemma is false for
some error function of the claimed form; the defect is that the printed
proof does not deliver it.

For orientation, the algebraic identity
$m\,\eps^{1/2} = (m^2\eps)^{1/2}$ shows what a derivation of the claimed
$\eps$-dependence would need: a chain that loses only square roots in the
$\eps$-type quantities \emph{after} the restricted-distribution
inflation would produce terms of the shape $(m^2\eps)^{1/2}$.  The
printed chain loses fourth roots, through two nested Cauchy--Schwarz
steps applied to already square-rooted quantities, and its
$md/q$-dependence enters as $m\,\delta_P^{1/4}$, which the claimed form
cannot absorb.

\section{Correction and alternative repairs}

The first direction below is the candidate used in the blueprint's first-route
discussion. The ten-session investigation does not establish full two-player
and consumer sufficiency, so no final source correction is adopted.

\paragraph{First-route error form (candidate correction).}
The operative error form is
$\delta_{\mathrm{combine}} = m \cdot \poly(\eps, md/q)$ --- concretely,
the first-route bound.  Following the proof-gap protocol, the source-labelled
statement keeps its printed form and the correction is carried by the proof,
by the chapter remark, and by a separate formalization-support node, rather
than being substituted silently into the source statement.  The only consumer
is the polynomial-pair lemma, whose imported error
$a(md)^a(\delta_{\mathrm{combine}}^b + q^{-b} + 2^{-bmd})$ absorbs
$m^{O(1)}$ prefactors by enlarging $a$ and shrinking $b$.  The blueprint proof
records this absorption and the corrected deficit reading of the second claim.
The unresolved obligations include explicit global-pair absorption and the
second-player X-removal and Z-deficit estimates. The first-player chain and
its integration into the actual conditional POVM are now proved.

\paragraph{Repairing the second route.}
Strengthen Property~2 of the sub-line lemma to the joint product law: in
each mixture component, the pairs $(\ell_X, x)$ and $(\ell_Z, z)$ are
independent with the stated restricted-distribution laws.  The sampling
procedure makes this plausible, since in each component the two pairs
are built from disjoint collections of the independent ingredients (the
two blocks of the uniform point, the direction blocks, and the fresh
auxiliary choices).  Open obligations: (i)~state and prove the
strengthened property for every component, including lines whose
direction blocks vanish and degenerate to singletons; (ii)~re-derive the
mixture decomposition of $\Lambda$ from the strengthened property;
(iii)~note that even then the route delivers only
$O(m^2\delta_P) = m^2 \cdot \poly(\eps, md/q)$, so this direction removes
the unsupported step of the printed proof but still requires the
error-form decision of the previous direction, with $m^2$ in place of
$m$.

\paragraph{Establishing the claimed form.}
Re-derive the chain so that every factor of $m$ enters jointly with
$\eps$ as $m^2\eps$, in the sense of the orientation identity above.
Open obligations: (i)~find a chaining that avoids one of the two nested
square-root losses, or prove restricted-distribution versions of the
underlying consistency lemmas directly at their inflated errors;
(ii)~handle the $md/q$-dependence, since a residual term
$m (md/q)^{c}$ is incompatible with the claimed form for the current
exponents.  No derivation along these lines is known here, and the
direction may be infeasible; in that case the corrected error form of
the first direction stands.

\section{Comparison with the blueprint}

The blueprint \cite{MIPStarREBlueprint} keeps the combined-lines lemma in
its source shape, error form $\poly(m^2\eps, md/q)$ included, and marks in
the proof that this form remains unproved: the proof derives the first route
only, records the established bound $m \cdot \poly(\eps, md/q)$, and uses
that bound in the sequel.  The
three claims are stated as separate lemmas with the corrected deficit
assertion for the second; the sub-line lemma states Property~2 at the
level of the two marginal laws and expressly disclaims the joint
conditional law of the two points, so the blueprint statement does not
license the second route; and the polynomial-pair lemma's proof absorbs
the $m$ prefactor into the constants of the imported soundness error.  A
chapter remark summarizes the discrepancy and cites this note.%
\footnote{The blueprint statements are \path{lem:qld-4-13},
\path{lem:qld-sublines}, \path{lem:claim-17-1}--\path{lem:claim-17-3},
\path{lem:qld-4-7}, and the remark
\path{rem:qld-4-13-source-defects}, all in
\path{blueprint/src/chapter/ch15_qpbt_combining.tex}.}

The source-labelled lemma links the carrier
\path{ExtendedLinesWitness} and the declaration
\path{exists_extendedLinesWitness}, which retains the printed error form
$\poly(m^2\eps,md/q)$ and whose proof is the tracked obligation.  The
established error form is linked from the separate formalization-support node
\path{lem:qld-4-13-established}, whose declaration
\path{exists_extendedLinesWitness_established} has the explicit form
$C m\,\poly(\eps,md/q)$ for a universal $C>0$; this agrees with the blueprint
convention that constants are absorbed into $\poly$.  Both proofs are open.

Neither declaration is source-shaped beyond its error form, and this note does
not claim otherwise.  Both return the carrier \path{ExtendedLinesWitness},
which replaces the source's lines in $\F_q^{2m+2}$ by the directly indexed
extended line space of
\path{docs/paper-gaps/qpbt_ld-dimension-divisibility.tex}, averages over the
corresponding directly indexed line-point law instead of over the line-point
distribution instantiated at dimension $2m+2$, and compares both sides of its
two consistency displays on the completed answer alphabet
$\F_q\cup\{\bot\}$, whose extra outcome collects the line polynomials that
admit no evaluation at the sampled point, in place of the source's answer
summation over $a\in\F_q$.  Neither blueprint node therefore carries a
statement-level formalization mark, and relating the two formulations is the
transport obligation of that note.

\section{The paired-line precursor: conditional collisions (issue 118)}

This section concerns the earlier assertion \texttt{lem:qld-xz-lines}, not a
new correction to the extended-line estimate discussed above. Its source is
\path{references/qpbt-paper/14_analysis_of_the_pauli_basis_test.tex},
lines 882--894, with the pasting argument at lines 943--963. It constructs
paired-line POVMs consistent with the already constructed joint points, with
error polynomial in $\eps$ and $\rho=md/q$, and with degree at most $d$ on
axis lines. Issue 118 records that the former unproved internal Lean helper
allowed an arbitrary supplied joint-point error $t$, while promising an
output error independent of $t$. This stronger quantification does not
follow from the printed argument. The internal interface now fixes a
polynomially controlled point-error function before choosing the line-error
function, as proved below. No counterexample to the printed existence
theorem is known here.

There is also a concrete obstruction to applying the conditional collision
hypothesis literally to every coefficient-polynomial fiber in Lean. For a
zero direction, every affine parameter denotes the base point. The completed
evaluation is defined only when all these parameters give the same value.
Consequently the distinct polynomials $T$ and $T+1$ both have completed value
$\bot$. If $1\le md<q$, they collide with probability one, exceeding $md/q$.
This is a counterexample to the unrestricted conditional-collision assertion,
not to the source game or to paired-line consistency. The source sampler in
\texttt{def:line-point-dist},
\path{references/qpbt-paper/08_classical_and_quantum_low_degree_tests.tex},
lines 274--287, and the evaluation convention are left unchanged.

\paragraph{Proved distributional estimates.}
Let $D$ denote that unchanged line-point law and let $E$ be the event that
the sampled direction is nonzero. On every line fiber in $E$, a uniform
affine parameter gives the correct conditional point law. More precisely,
for every real function $H$ on line-point pairs,
\[
 \mathbb E_{(\ell,x)\sim D}H(\ell,x)
 =\mathbb E_{(\ell,x)\sim D}\mathbb E_{s\in\mathbb F_q}
     H(\ell,\ell_0+s\operatorname{dir}(\ell)).
\]
The identity itself holds also for zero directions. Its proof preserves the
canonical representative jointly with the point: the representative of
$\ell_0+s\operatorname{dir}(\ell)$ is $\ell_0$, and resampling a uniform
point followed by a uniform parameter preserves the uniform point law.
The block decomposition of the conditionally linear sampler gives this
identity separately for axis and diagonal lines and then for their mixture.

For distinct coefficient polynomials $f,g$ of degree at most $b$, their
difference is a nonzero polynomial of degree at most $b$. It has at most
$b$ roots. In Lean, the coefficient tuple is mapped injectively to Mathlib's
polynomial type by \path{Polynomial.ofFn}; equality on more than $b$ field
elements would imply polynomial equality by
\path{Polynomial.eq_of_natDegree_lt_card_of_eval_eq}, and hence equality of
the original tuples. This explains the library replacement for the
Schwartz--Zippel invocation at source line 955. Together with resampling it
proves, for every nonnegative function $w$ of the line alone,
\[
 \mathbb E_D[\mathbf1_E w(\ell)
  \mathbf1_{\mathrm{eval}_x(f)=\mathrm{eval}_x(g)}]
 \le \frac bq\mathbb E_D[\mathbf1_E w(\ell)].
\]
Taking $w$ to be a line-fiber indicator proves the required collision ratio
on each nondegenerate positive-mass line fiber. The claim remains true when
$b\ge q$; in that case the bound need not be smaller than one.

Axis directions never vanish. A zero prefix-projected diagonal direction
forces the last raw direction coordinate to vanish, and this coordinate is
uniform, independently of the seed. Thus
\[
 D(E^c)\le\frac1{2q},\qquad D(E)\ge\frac34.
\]
The latter bound uses only $q\ge2$. On the product of two line-point laws,
discarding only zero $X$ directions has the same bound. This suffices for the
one-sided pasting orientation: the constructed POVM is $XZX$, so $G_1=Z$
and $G_2=X$, with the two answer coordinates exchanged. There is no need to
exclude degenerate $Z$ directions for the collision hypothesis.

Finally, if an event $F$ has mass $r>0$, the exact restricted-average identity
and the bound $h\le1$ give
\[
 \mathbb E_\mu h\le r\mathbb E_{\mu\mid F}h+\mu(F^c).
\]
For $h\ge0$, conditioning also gives
$\mathbb E_{\mu\mid F}h\le\mathbb E_\mu h/r$.
These identities retain the discarded mass rather than changing the game.
All the distributional assertions in this paragraph are proved in
\path{MIPStarRE/QPBT/Combining/Lines.lean} and linked from the separate
blueprint node \texttt{lem:combined-line-nondegenerate-collisions}.

\paragraph{Conditioned product law and actual marginals.}
The product collision predicate is now proved, not assumed. With
$s_X=(\ell_X,x)$ and $s_Z=(\ell_Z,z)$, normalize the restriction of
$D\otimes D$ to $\operatorname{dir}(\ell_X)\ne0$ and push it forward by
\[
 (s_X,s_Z)\longmapsto (((\ell_X,\ell_Z),s_Z),s_X).
\]
The first coordinate is the common line-pair question; the next two are the
Z and X evaluation questions, respectively. Retaining the line descriptors
inside the latter questions allows the completed evaluations to be ordinary
functions of their own question. On the support these descriptors agree
with the common question. No positive-mass fiber permits inconsistent
descriptors. Fixing the common question and the Z sample gives a
nonnegative weight depending only on the X line. The weighted root bound,
followed by averaging over Z, proves the collision estimate. Multiplying
through by the positive retained mass removes the normalization factor.
Thus the exact predicate required by one-sided pasting holds with $b/q$;
Z directions, including zero directions, remain in the law.

For the supplied point witness, the two completed marginals with answer
order $(Z,X)$ have now been transported to this law. Their defects are at
most $[8t+C(\eps+\sqrt\eps)]/r$. The line families depend only on the
common question, as required by pasting. Joint self-consistency is also
transported with bound $t/r$, in the common $(X,Z)$ order. These statements
are for every directed opposite placement, with no symmetry assumption.
The pointwise defect of any two complete measurements on opposite registers
equals one minus a nonnegative diagonal overlap. It is therefore at most
one. Consequently their original and conditioned defects satisfy
$d\le r d_E+1/(2q)$. This inequality has been specialized to the actual
$XZX$ line POVM and the supplied completed points; it does not assert a
bound on their conditional defect. These results are linked by the separate
auxiliary blueprint node \texttt{lem:combined-line-conditioned-product}.

\paragraph{Unequal local spaces and the unused EPR pair.}
The local spaces need not have equal dimensions. Let $H_L,H_R$ denote the
two sides of a bipartition and include $H_L\otimes H_R$ isometrically in
$(H_L\oplus H_R)\otimes(H_L\oplus H_R)$, using only its left-right
summand. Write $U$ for this inclusion. For arbitrary operators on the
indicated spaces,
\[
 \langle U\psi,(A\oplus B)\otimes(C\oplus D)U\psi\rangle
   =\langle\psi,A\otimes D\,\psi\rangle.
\]
This is not symmetrization: there is no exchanged copy of the state and no
factor of $1/2$. Its proof expands the finite coordinate sums; every term
outside the left-right sector is zero. In particular, the inclusion
preserves norms and consistency defects exactly.

Take the block sum of the supplied projective joint families on $H_L$ and
$H_R$. On the unused left block, complete each right line measurement by a
deterministic measurement. These outcomes exist because a normalized state
has nonzero right space, the question law has nonempty support, and a
complete measurement on a nonzero space has an outcome. These deterministic
blocks never replace the supplied measurements: the displayed identity
shows that they contribute nothing to any comparison. Direct sums commute
with outcome postprocessing and ordered products and preserve projectivity.
Applying the existing one-sided pasting theorem on the common carrier
therefore proves its heterogeneous version, with the same error function,
the same two forward marginal errors, and the self-consistency comparison
between the two supplied joint families. No reverse marginal comparison or
equality of those families is required. Both supplied joint families are
projective, as the point witness already guarantees.

For the expanded state, all four directed bipartitions have now been
constructed: $AA'\mid BA''(B'B'')$, $BA''\mid AA'(B'B'')$,
$BB'\mid AB''(A'A'')$, and $AB''\mid BB'(A'A'')$.
The unused EPR pair remains in the right space; it is not traced out or
replaced by an independence assumption. Coordinate reindexing preserves the
placed consistency defects. Adjoining the identity of this unused factor
commutes with evaluated pasting. Exchanging the two completed answers
preserves the supplied joint self-consistency bound $t/r$, and evaluating
the actual line POVM in the order $(Z,X)$ gives exactly the pasting output
with $G_1=Z,G_2=X$. These are proved representation identities, valid even
for zero directions; the collision estimate still uses only the previously
specified nonzero-X restriction. They are recorded in the separate
blueprint auxiliary \texttt{lem:combined-line-bipartite-transport}.

The coordinate identities hold also for empty finite carriers, when every
state is zero; normalization excludes this case in the pasting application.
Unequal dimensions introduce no normalization loss. The answer-exchange and
unused-factor identities require neither injective evaluation nor a product
state. Thus they introduce no exception at degree zero, at $q=2$, or on a
zero Z direction. These are symbolic boundary checks covered by the proofs,
not a claim to have enumerated all finite strategies.

\paragraph{Application to the actual conditional families.}
The supplied joint-point marginals already have proved bounds
$u=8t+C(\eps+\sqrt\eps)$ for a universal $C\ge1$; joint self-consistency
has bound $t$. Conditioning on nondegenerate $X$ lines therefore retains
both errors, with the marginal bound $u/r$. These are now the actual inputs
to heterogeneous pasting: the right line families are $Z$ then $X$, and
the two joint families are the supplied completed points with answers
exchanged. The point witness itself implies $t\ge0$, so $t/r\le u/r$;
no new error hypothesis is needed. Pulling the conditioned question law back
to the original line-point samples makes the common descriptor and its
local duplicate identical. This is an identity of averages under the defined
map, not an assumption that all ambient questions have consistent descriptors.
The four bipartitions and the evaluated sandwich identity then give
\[
 d_E\le f(\rho,u/r),\qquad
 d\le r f(\rho,u/r)+\frac1{2q},
\]
where $f$ is the same universal additive polynomial error function supplied
by pasting. The second bound exchanges both answers back to $(X,Z)$ and
applies the proved restoration inequality. In particular neither $r$ nor
the discarded contribution is dropped. These two inequalities are proved in
\leanid{MIPStarRE.QPBT.exists\_combinedLine\_conditioned\_defect\_le} and
\leanid{MIPStarRE.QPBT.exists\_combinedLine\_restored\_defect\_le}, and linked
by the separate blueprint auxiliary
\texttt{lem:combined-line-actual-pasting-bound}.

The construction applies to unequal local dimensions, all four directed
placements, and differing supplied joint families. It retains zero $Z$
directions. For $q=2$ the retained mass is still positive. The two
polynomials $T,T+1$ on a zero $X$ direction no longer satisfy the retained
event; their original-law contribution is bounded by the restored term.
Positive supplied error $t$ is not silently set to zero when $\eps=0$.
These are symbolic counterexample and boundary checks, not an exhaustive
numerical search over quantum strategies.

\paragraph{Remaining quantitative specialization.}
For the explicit pasting estimate $D_0(\rho^{1/4}+(u/r)^{1/8})$, the
corresponding scalar calculation is
\[
 rD_0\bigl(\rho^{1/4}+(u/r)^{1/8}\bigr)+(1-r)
 \le D_0\bigl(\rho^{1/4}+u^{1/8}\bigr)+\frac1{2q}.
\]
Here $0<r\le1$ accounts for both powers of $r$; no error term is omitted.
For the constructed points, $t=K\eps^{1/8}$. On $0\le\eps,\rho\le1$,
$u\le(8K+2C)\eps^{1/8}$, while $1/(2q)\le\rho/2$ since $md\ge1$.
Thus this calculation has the additive polynomial form
$O(\rho^{1/4}+\eps^{1/64})$. Larger parameters can use the unit bound for
normalized POVM consistency. The new Lean inequalities expose a general
additive polynomial function, not these particular exponents. More generally,
if $f(x,y)\le A(x^a+y^b)$ with $a,b>0$, the mass bound $r\ge3/4$
gives
\[
 r f(\rho,u/r)\le
 A\rho^a+A\max\{1,(4/3)^{b-1}\}u^b.
\]
On the unit square this gives an additive bound with exponents
$b/8$ in $\eps$ and $\min\{a,1\}$ in $\rho$. This explains why the
existential pasting interface suffices mathematically without assuming
monotonicity of its unspecified function. Formal scalar composition and
assembly into the existing line witness, including exchange of the opposite
placements, are now proved as described in the following subsection. The
extended-line construction in \path{Combining/Apply.lean} remains separate.
The equal-carrier interface is no longer a
representation obstruction: the proved isometric inclusion eliminates it
without assuming equal player dimensions or equal supplied point families.
The source existence theorem is still not disproved, and neither final
mathematical sufficiency nor minimality of a source correction is established.

The affected source uses are the restricted averages at lines 1058--1061,
extended-line construction and degree at 1118--1132, and consistency at
1241--1246; global pairs at 1278--1282 and extraction at 1423, 1465--1492,
1584, 1627, and 1671 depend on these. The complete blueprint-descendant
inventory is preserved in
\path{audits/2026-09-06-issue118-point-error-dependency.md}.
Neither the new auxiliary node nor the algebraic orientation identity marks
any of those source results as proved. Claims 17-2 and 17-3, the public
existence statements, and all witness fields remain unchanged.
The product collisions, conditioned supplied marginals, heterogeneous
operator application, and restored consistency concern the actual POVMs.
The following specialization completes the paired-line witness, but not
the affected extended-line construction. No final correction or minimality
claim for the extended-line theorem is adopted, and no change to a
mathematical definition or game is requested.

\subsection{Polynomial specialization and the first-route overlap}

Fix a function $t$ with $0\le t(\eps)\le K\eps^\alpha$ for all
$\eps\ge0$, where $K\ge1$ and $\alpha>0$. The corrected internal
quantifier order is
\[
  \forall t\text{ with this bound},\quad
  \exists g=\poly(\eps,\rho),\quad
  \forall P,\eps,S,Q\text{ with point error }t(\eps),\quad
  d(T,Q)\le g(\eps,\rho).
\]
Here $T$ is still the actual $XZX$ measurement, and the conclusion includes
axis-degree support and all four directed comparisons. The earlier
restored bound remains available for arbitrary numerical $t$; this
specialized interface does not discard its dependence on $t$.

Write $f(x,y)\le A(x^a+y^b)$ with $A\ge1$ and $a,b>0$. Set
\[
 \gamma=\min\{\alpha,1/2\},\quad \beta=\min\{a,1\},\quad
 D=8K+2C,\quad B=A(1+(2D)^b)+1.
\]
For $0\le\eps,\rho\le1$, the elementary ordering of powers gives
$u=8t(\eps)+C(\eps+\sqrt\eps)\le D\eps^\gamma$.
The retained mass satisfies $1/2\le r\le1$, so $u/r\le2D\eps^\gamma$.
One applies monotonicity to the explicit power bound for $f$, not to $f$
itself. Since $md\ge1$, the restored error obeys
\[
 \min\left\{1,r f(\rho,u/r)+\frac1{2q}\right\}
 \le B\bigl(\eps^{\gamma b}+\rho^\beta\bigr).
\]
If either parameter exceeds one, the right-hand side is at least one,
and the independent unit bound for normalized POVM consistency applies.
The minimum is used only in this scalar estimate; it does not choose a
different measurement or modify the question distribution. Taking
$t(\eps)=K\eps^{1/8}$, as supplied by the point construction, specializes
$\gamma$ to $1/8$. With the explicit pasting powers $a=1/4$, $b=1/8$,
this gives $O(\eps^{1/64}+\rho^{1/4})$.

This argument proves
\leanid{MIPStarRE.QPBT.exists\_conditioned\_polynomial\_bound}.
Opposite-placement operators commute, and exchanging the two outcome
summations therefore preserves their consistency defect. This observation,
formalized by
\leanid{MIPStarRE.QPBT.consistencyDefect\_opposite\_symm}, supplies the
orientation required by the existing line witness without any assumption
on equality of dimensions, equality of point families, or symmetry of the
state. Consequently
\leanid{MIPStarRE.QPBT.exists\_combinedLinesWitness\_ofPointsWitness}
and \leanid{MIPStarRE.QPBT.exists\_combinedLinesWitness} now have proofs.
The latter's public statement is unchanged: it supplies the polynomial
point function and point witness internally.

For the first extended-line route, let $J$ denote the completed joint
overlap on the subline law. The triangle inequality applied to the existing
Claims 17-1, 17-2, and 17-3 gives a universal $C_*>0$ such that
\[
 |J-1|\le C_*\bigl(t^{1/2}+\sqrt m\,(p^{1/4}+t^{1/4}+\eps^{1/4})\bigr),
\]
where $p$ is the supplied paired-line error. This is proved by
\leanid{MIPStarRE.QPBT.subline\_joint\_overlap\_near\_one}, without
altering any of the three claims. If $p\le A(\eps^a+\rho^b)$ and
$t\le K\eps^\alpha$, fourth-root subadditivity and $\sqrt m\le m$
bound the expression, for $\eps\le1$, by a constant times
$m(\eps^\theta+\rho^{b/4})$, where
$\theta=\min\{\alpha/4,a/4,1/4\}>0$. For $\eps>1$ that polynomial
can be chosen at least one. This proves the scalar assertion
\leanid{MIPStarRE.QPBT.exists\_combining\_polynomial\_bound}, whose left
side is the minimum of one and the displayed error expression. It does
not assert that the extended-line measurement has already been constructed.

The two internal extended-line helpers now also fix the polynomial
point-error function before their output error functions are chosen.
Their existential wrappers still produce that function internally and
have unchanged statements. Three existing obligations remain in
\path{Combining/Apply.lean}: the printed-error construction, the
first-route extended-line construction, and the global-pair construction.
For the first-route obligation, the conditional average, axis-degree support,
and point-effect comparison are now proved as described below. One must still
integrate the placed overlaps and establish the second-player first-route
estimate. The source-question and field-answer transport obligations in
the earlier section are unchanged. In particular neither the new scalar
estimate nor the overlap estimate establishes the printed stronger error
$\poly(m^2\eps,md/q)$.

The counterexample audit covers $\eps=0$, the scalar boundary $\rho=0$,
$q=2$, arbitrary positive polynomial exponents, large errors, zero $Z$
directions, unequal player dimensions, and unequal supplied point families.
The zero-$X$ collision counterexample is excluded only from the proof law,
and its original-law contribution is restored. An arbitrary positive
supplied error at $\eps=0$ no longer satisfies the specialized polynomial
contract. These are symbolic arguments, not a numerical search over all
quantum strategies. Restricting the internal function to exactly
$K\eps^{1/8}$ would be stronger than the proved sufficient polynomial
contract. No minimality claim for the unresolved extended-line theorem is
made.

\section{The actual conditional extended-line construction}

The construction at source lines 1118--1132 can be carried out without
changing any game, line carrier, or evaluation convention. Let $D$ be a
subline witness and let $T$ be the supplied paired-line POVMs. Every
canonical direct line has positive marginal mass: its stored base is fixed
by the canonical projection, and its stored direction satisfies the prefix
condition in the diagonal case. Sampling that base, index, and direction
therefore produces the line with positive probability. This argument
includes diagonal lines with zero direction. The marginal property of $D$
then makes each supported fiber have positive mass.

Normalize that fiber, choose affine coordinates whose compatibility is
proved by the subline witness, and apply the source polynomial combination
to the actual outcomes of $T$. The resulting conditional average is a
complete POVM: positive scalar combinations preserve positivity, and
interchanging the two finite sums gives the identity operator. No
replacement POVM, default answer, or zero-mass branch is used. Finite
disintegration cancels the positive fiber mass and recovers the original
law $D$. These assertions are proved by
\leanid{MIPStarRE.QPBT.SubLineWitness.extendedMeasurement},
\leanid{MIPStarRE.QPBT.SubLineWitness.conditionalLaw\_isProbability}, and
\leanid{MIPStarRE.QPBT.SubLineWitness.avgOver\_conditionalLaw}.

For axis lines, there are two cases. If the direction is not a scalar
coordinate, both scalar weights are constant; affine substitution preserves
the two source degree bounds. If the direction is a scalar coordinate,
both projected directions vanish. Compatibility at parameters zero and one,
together with the nonzero source axis directions, forces both projected
parameter slopes to vanish. The two composed source polynomials are then
constant, and the combined polynomial has degree at most one, hence at most
$d$. Coefficients above $d$ vanish algebraically, even when $d\ge q$;
this argument does not infer polynomial equality from equality of functions
on the field. The zero effects of invalid source outcomes pass through
postprocessing and averaging. This proves
\leanid{MIPStarRE.QPBT.SubLineWitness.extendedMeasurement\_axis\_degree}.

At a point $u=(x,z,\alpha,\beta)$, suppose both completed source evaluations
are determined values $a,b$. Every parameter representing $u$ also presents
the projected points, so the combination formula gives $\alpha a+\beta b$
for every such parameter. Thus the combined completed evaluation is the
corresponding field answer, including the zero-direction case. For any
point-pair POVM, its effect at $(a,b)$ is a positive summand of its linear
postprocessing at $\alpha a+\beta b$. If either source evaluation is
undefined, the completed point-pair effect is zero instead. This gives
\leanid{MIPStarRE.QPBT.SubLineWitness.completedPoint\_effect\_le}
on either player space, without assuming equal dimensions or a symmetric
strategy. The earlier collision counterexample $T,T+1$ on a zero direction
does not violate this assertion: both source evaluations are undefined,
and the compared point effect is zero.

These kernel-checked assertions establish the actual POVM and its degree
support, but not the two quantitative consistency fields of the extended-line
witness. The following integration results further reduce the remaining work.
All three existing holes in
\path{MIPStarRE/QPBT/Combining/Apply.lean} remain explicitly visible.
Neither the printed stronger error, source-question and field-answer
transport, nor global-pair sufficiency is certified by this increment.
The complete consumer inventory remains in
\path{audits/2026-09-06-issue118-point-error-dependency.md}.

\section{Integration and final ordinary-budget disposition}

Write $J_{p_1,p_2}$ for the actual paired-line versus completed point-pair
overlap on $D$ with a uniform affine parameter. Uniform resampling preserves
the direct line descriptor jointly with the sampled point, including zero
directions. Expand postprocessing and the conditional average inside the
state overlap. Positive operators on opposite registers preserve the
point-effect inequality proved above. Exchanging the finite averages and
disintegrating $D$ gives, for every directed opposite placement,
\[
 \operatorname{Defect}_{p_1,p_2}(Q^{\rm line},Q^{\rm point})
 \le 1-J_{p_1,p_2}.
\]
This is proved by
\leanid{MIPStarRE.QPBT.SubLineWitness.extended\_consistencyDefect\_le};
it makes no assertion about a joint product law conditional on source lines.

For any placements, the difference between $J$ and the ordered
$Z$-then-$X$ point-product overlap is at most $\sqrt{4t}$. The completed
operator squared distance is at most $4t$, and its subline average equals
its uniform joint point average; Cauchy--Schwarz then gives the estimate.
This proves the second-player counterpart of the first replacement step,
\leanid{MIPStarRE.QPBT.SubLineWitness.paired\_ordered\_overlap\_gap\_le}.
Combining the integration inequality with the previously proved first-player
three-claim estimate proves
\leanid{MIPStarRE.QPBT.extended\_first\_player\_deficit\_le} for the actual
conditional POVM, with bound
$C(t^{1/2}+\sqrt m(p^{1/4}+t^{1/4}+\eps^{1/4}))$.

The remaining second-player work is specifically the X-removal and
Z-deficit estimates at placements $(\mathsf{BB}',\mathsf{AB}'')$.
The existing one-point deficit helpers are specialized to the first-player
orientation; neither commutation of opposite operators nor the ordered-product
replacement alone transports those estimates. Once they are proved, the
existing scalar absorption and degree-support construction can assemble the
first-route witness. This still would not establish the stronger printed
error or certify global-pair and extraction sufficiency.

Attempt 10 is the final ordinary attempt. The cumulative completed time
carried into it is 17558 seconds, including the dispatcher's final
2406-second duration for attempt 9; the original anchor remains
September 5, 2026, 19:24 UTC. The complete owner-escalation dossier, including
all ten attempts, the exact remaining declarations, and the source and
blueprint consumer inventory, is recorded in
\path{audits/2026-09-06-issue118-point-error-dependency.md}.
No additional attempt, source correction adoption, definition change, or
publication is authorized by this record.

\section{Conclusion}

There is a proof-level gap in the cited lemma, together with an
intermediate inconsistency: the printed proof asserts two different
bounds for the same consistency defect, the first-route bound
$O(m(\eps^{1/4} + \delta_P^{1/4} + \delta_Q^{1/4} +
\delta_{\mathrm{Line}}^{1/2}))$ and the second-route bound
$O(m^2\delta_P)$, the second resting on a joint-law decomposition that
the sub-line lemma as stated does not supply; and the error form claimed
in the lemma statement, $\poly(m^2\eps, md/q)$, follows from neither.
The first-route error bookkeeping gives $m \cdot \poly(\eps, md/q)$,
and the global-pair error expression can absorb polynomial-in-$m$
prefactors. These mathematical observations do not discharge the remaining
two-player construction and source-transport obligations. The blueprint therefore
retains the printed statement with the missing step explicitly tracked, uses
the corrected first-route form in the proof and in the sequel, and cites this
note; strengthening the sub-line lemma to license the second route is
optional, and no route to the printed form is currently known.  We do not
assert that the printed form is false.

\bibliographystyle{alpha}
\bibliography{references}

\end{document}
